% Bewertungspaket zum Gesamttest «Dreiecke» — Quelle des PDFs (python3 scripts/build-lp-pdf.py leitprogramme/dreiecke/).
% Ganze Punkte je Teilschritt; Werte mit python3 nachgerechnet (scripts/lp/dreiecke/zahlen.py, 08.10.2026, Fehlerfälle
% eingeschlossen). Aufbau wie trigonometrische-berechnungen/bewertungspaket.tex.
\documentclass[11pt]{article}
\usepackage{../lp-druck}
\renewcommand{\lpname}{Dreiecke}
\renewcommand{\lpbereich}{Grundlagenfach}
\renewcommand{\lpteilgebiet}{5.2a}
\renewcommand{\lpfuss}{Bewertungspaket}
\newcolumntype{P}{>{\centering\arraybackslash}p{10mm}}
\newenvironment{raster}{\par\smallskip\noindent\tabularx{\linewidth}{@{}>{\raggedright\arraybackslash}p{38mm}P X@{}}\toprule
  \small\textsc{Teilschritt} & \small\textsc{P} & \small\textsc{Musterlösung}\\\midrule}{\bottomrule\endtabularx\par}
\newcommand{\fehler}[1]{\par\smallskip{\small\textcolor{lprot}{\bfseries Typische Fehler:} #1}}
\newcommand{\E}{\,{\footnotesize\bfseries\color{lpgruen}(E)}}
\newcommand{\A}{\,{\footnotesize\bfseries\color{lpblau}(A)}}
\begin{document}
\lpkopf{Bewertungspaket zum Gesamttest}{}

\begin{lpachtung}{Erst lösen, dann öffnen}
Dieses Paket enthält die Musterlösung. Wer es vor dem Lösen liest, misst am Ende nur, wie gut das Abschreiben gelingt.
\end{lpachtung}

\section*{1 · So gehst du vor}
\begin{enumerate}[leftmargin=*,itemsep=0pt]
\item \textbf{Gesamttest lösen} — vollständig, mit Skizze und Rechenweg, ohne dieses Paket; der Taschenrechner ist erlaubt.
\item \textbf{Fotografieren oder scannen}, gut lesbar, eine Seite pro Bild. \textbf{Kein Name} auf den Bildern.
  Handyfotos tragen oft unsichtbar Standort, Zeit und Gerät mit (Metadaten). Lade darum lieber einen Scan oder ein
  Bildschirmfoto des Fotos hoch, oder entferne vorher die Standortangaben.
\item \textbf{Bewerten lassen.} Ob du dafür eine KI nutzt und welche, entscheidest du selbst und in eigener Verantwortung —
  je nachdem, was dir aufgrund deines Alters und deiner persönlichen Situation erlaubt ist (Altersgrenzen und
  Nutzungsbedingungen des Anbieters, allenfalls Einverständnis deiner Eltern). Lade nur deine Lösung und dieses PDF hoch,
  keine persönlichen Daten, und füge den Auftrag aus Abschnitt~2 ein. Ohne KI kannst du dich mit Abschnitt~3 selbst bewerten.
\item \textbf{Rückmeldung prüfen.} Stimmt, was die KI aus deiner Handschrift gelesen hat? Eine KI kann sich irren —
  im Zweifel gilt das Raster in Abschnitt~3.
\item \textbf{Gezielt wiederholen} nach Abschnitt~4.
\end{enumerate}

\needspace{8\baselineskip}
\section*{2 · Auftrag an die KI}
\begin{tcolorbox}[breakable,colback=lplinie!25,colframe=lplinie,boxrule=0.5pt,arc=1mm,fontupper=\small,left=2mm,right=2mm,top=1.5mm,bottom=1.5mm]
Du bist Korrektorin oder Korrektor für Mathematik an einer Schweizer Berufsmaturitätsschule. Im Anhang findest du (1)~das Bewertungspaket zum Gesamttest «Dreiecke» mit Musterlösung und Punkteraster und (2)~Fotos meiner handschriftlichen Lösung.\par\smallskip
Geh so vor:\par
1. Schreib zu jeder Aufgabe G1 bis G7 zuerst ab, was du in meiner Lösung liest — Rechenweg und Ergebnis. Was du nicht sicher lesen kannst, markierst du mit [unsicher]. Nicht raten.\par
2. Vergib die Punkte genau nach dem Raster im Bewertungspaket (Abschnitt 3), ganze Punkte je Zeile. Ziehe einen Fehler nur einmal ab: Rechne ich mit einem falschen Zwischenergebnis richtig weiter, gibt es die Folgepunkte. Ausnahme: Zeilen mit \textbf{(E)} gibt es nur für das richtige Ergebnis, nie als Folgepunkt. Die Zeilen «Typische Fehler» sagen, welche Rasterzeile bei diesem Fehler 0 Punkte gibt — sie sind kein zusätzlicher Abzug.\par
3. Andere richtige Lösungswege zählen voll, ebenso gleichwertige Schreibweisen: Dezimalpunkt oder Dezimalkomma, exakte Werte wie $\sqrt{160}$ statt Dezimalzahlen, anders, aber vernünftig gerundete Ergebnisse ($\pm 0.01$ oder durch gerundete Zwischenwerte erklärbar). \textbf{Einheiten:} Fehlt bei einem Endergebnis die Einheit oder ist sie falsch ($\text{cm}$ statt $\text{cm}^2$), gibt die (E)-Zeile dieses Ergebnisses 0 Punkte — aber nur beim ersten Mal im ganzen Gesamttest; weitere Einheitenfehler kosten nichts mehr. Zeilen mit \textbf{(A)} verlangen nur Erkennen oder Ablesen, diese Punkte gibt es auch ohne Rechenweg. Zeilen mit \textbf{(E)} sind Ergebnispunkte: Das Ergebnis muss stimmen \emph{und} aus einem erkennbaren Weg hervorgehen. Alle übrigen Zeilen bewerten den Weg selbst.\par
4. Begründe jeden Abzug in einem Satz und nenne die Fehlerart (zum Beispiel «schräge Seite als Höhe» oder «Hypotenuse als Kathete eingesetzt»). Schreib die Musterlösung nicht ab — zeig mir, wo mein Fehler liegt.\par
5. Gib zum Schluss eine Tabelle «Aufgabe | Punkte | Begründung», die Gesamtpunktzahl (von 25) und — nach der Selbsteinschätzung in Abschnitt 4 — welches Kapitel des Leitprogramms ich wiederholen soll. Keine Note.\par
6. Fehlt ein Foto oder ist eine Aufgabe unlesbar, sag es, statt sie zu bewerten.
\end{tcolorbox}

\needspace{8\baselineskip}
\section*{3 · Musterlösung und Punkteraster}
\textbf{Allgemein:} Ganze Punkte je Zeile. Ein Fehler wird nur einmal abgezogen (Folgefehler: wer mit einem falschen
Zwischenergebnis richtig weiterrechnet, bekommt die Wegpunkte der folgenden Zeilen, aber keine (E)-Punkte). In der Spalte P:
\textbf{(A)}~= erkennen oder ablesen, zählt auch ohne Rechenweg · \textbf{(E)}~= Ergebnispunkt, Ergebnis richtig
\emph{und} Weg erkennbar · ohne Zeichen = die Zeile bewertet den Weg. Taschenrechner erlaubt; Längen und Winkel auf zwei
Dezimalen, Toleranz $\pm 0.01$ (oder durch gerundete Zwischenwerte erklärbar). \textbf{Einheiten:} Der erste Einheitenfehler im
ganzen Test kostet die (E)-Zeile dieses Ergebnisses, weitere nichts.
«Typische Fehler» nennt, welche Zeile 0 Punkte gibt, und die Punkte, die bleiben.
\textbf{Teilrichtige Zeilen:} Verlangt eine Zeile mehrere Angaben und stimmt nur ein Teil, gibt die Zeile 0 Punkte.

\begin{lpaufgabe}{G1}{Winkel aus dem Aussenwinkel}{4 P}
\begin{raster}
(a) Ansatz & 1 & Aussenwinkelsatz: $\gamma' = \alpha + \beta = \alpha + 2\alpha = 3\alpha = 126^\circ$ (gleichwertig: $\gamma = 180^\circ - 126^\circ = 54^\circ$, dann $3\alpha = 126^\circ$)\\
(a) $\alpha$, $\beta$ & 1\E & $\alpha = 42^\circ$, $\beta = 84^\circ$. Beide nötig.\\
(a) $\gamma$ & 1 & $\gamma = 180^\circ - 126^\circ = 54^\circ$ (oder $180^\circ - 42^\circ - 84^\circ$). Folgepunkt: richtig aus den eigenen Winkeln.\\
(b) Art & 1\A & Spitzwinklig, weil alle drei Winkel unter $90^\circ$ liegen (der grösste ist $84^\circ$); nicht gleichschenklig, weil keine zwei Winkel gleich sind. Beide Aussagen mit Grund nötig. Folgepunkt: richtig aus den eigenen Winkeln.\\
\end{raster}
\fehler{Aussenwinkel als Innenwinkel $\gamma = 126^\circ$ genommen ($3\alpha = 54^\circ$, $\alpha = 18^\circ$, $\beta = 36^\circ$): Ansatz und Ergebniszeile 0; $\gamma$-Zeile 0; (b) als Folgepunkt («stumpfwinklig, weil $126^\circ$») $\to$ 1 von 4\,P.
$\beta = \alpha$ statt $2\alpha$ ($\alpha = \beta = 63^\circ$): Ansatz und Ergebniszeile 0, $\gamma = 54^\circ$ zählt; (b) als Folgepunkt (gleichschenklig) $\to$ 2 von 4\,P.}
\end{lpaufgabe}

\begin{lpaufgabe}{G2}{Brunnen und Bank im Park}{4 P}
\begin{raster}
(a) Brunnen & 1\A & Die Mittelsenkrechten der Seiten (zwei genügen); der Punkt ist der Umkreismittelpunkt $M_U$. Beides nötig.\\
(b) Lage & 1 & Nein: Das Dreieck ist stumpfwinklig ($105^\circ > 90^\circ$), darum liegt $M_U$ ausserhalb. Ohne Grund: 0.\\
(c) Bank & 1\A & Die Winkelhalbierenden (zwei genügen); der Punkt ist der Inkreismittelpunkt $M_I$. Beides nötig.\\
(d) Lage & 1 & Nein: $M_I$ liegt immer im Dreieck (die Winkelhalbierenden verlaufen im Innern). Ohne Grund: 0.\\
\end{raster}
\fehler{Brunnen und Bank vertauscht (Winkelhalbierende für die Ecken, Mittelsenkrechte für die Wege): (a) und (c) 0; (b) und (d) zählen, wenn die Begründung zum genannten Punkt stimmt $\to$ höchstens 2 von 4\,P.
Höhen statt Mittelsenkrechten: (a) 0.}
\end{lpaufgabe}

\begin{lpaufgabe}{G3}{Höhe und Winkelhalbierende}{4 P}
\begin{raster}
$\gamma$ & 1 & $\gamma = 180^\circ - 74^\circ - 46^\circ = 60^\circ$\\
Winkel an der Höhe & 1 & Im rechtwinkligen Dreieck $AFC$: $\angle ACF = 90^\circ - 74^\circ = 16^\circ$ (gleichwertig über $B$: $\angle BCF = 90^\circ - 46^\circ = 44^\circ$)\\
Halber Winkel & 1 & $\angle ACD = \tfrac{\gamma}{2} = 30^\circ$ (bzw. $\angle BCD = 30^\circ$)\\
Ergebnis & 1\E & $30^\circ - 16^\circ = 14^\circ$ (bzw. $44^\circ - 30^\circ = 14^\circ$)\\
\end{raster}
\fehler{Höhe als Winkelhalbierende behandelt oder $\angle ACF = \alpha$: Zeile «Winkel an der Höhe» 0, das Ergebnis 0; $\gamma$ und der halbe Winkel zählen $\to$ 2 von 4\,P.
Summe statt Differenz ($30^\circ + 16^\circ = 46^\circ$): nur die Ergebniszeile 0 $\to$ 3 von 4\,P.}
\end{lpaufgabe}

\begin{lpaufgabe}{G4}{Eine fremde Lösung}{4 P}
\begin{raster}
(a) Fehler & 1 & Mia nimmt die schräge Seite $a = 5\,\text{cm}$ als Höhe. Zur Grundseite $c$ gehört die Höhe $h_c = 4\,\text{cm}$, die senkrecht auf der Geraden $AB$ steht (ihr Fusspunkt liegt ausserhalb).\\
(a) Fläche & 1\E & $A = \tfrac{1}{2} \cdot 9 \cdot 4 = 18\,\text{cm}^2$\\
(b) Seite $b$ & 1 & $AC$ ist die Hypotenuse des rechtwinkligen Dreiecks $AFC$ mit $\overline{AF} = 9 + 3 = 12\,\text{cm}$ und $h_c = 4\,\text{cm}$: $b = \sqrt{12^2 + 4^2} = \sqrt{160} \approx 12.65\,\text{cm}$\\
(b) Umfang & 1\E & $U = 9 + 5 + 12.65 \approx 26.65\,\text{cm}$. Folgepunkt nur für die Wegzeile, nicht für (E).\\
\end{raster}
\fehler{$\overline{AF} = 9\,\text{cm}$ statt $12\,\text{cm}$ ($b = \sqrt{97} \approx 9.85$, $U \approx 23.85$): Seite-$b$-Zeile 0, Umfang (E) 0 $\to$ 2 von 4\,P.
Umfang ohne $b$ oder mit $b = 9 + 3$ ($U = 26$): beide (b)-Zeilen 0.
Fehler genannt, aber Fläche mit $\tfrac{1}{2}$ vergessen ($36$): (a) Fläche 0.}
\end{lpaufgabe}

\begin{lpaufgabe}{G5}{Rechtwinkliges Dreieck in schräger Lage}{3 P}
\begin{raster}
(a) Hypotenuse & 1\A & $\overline{KM}$ — sie liegt dem rechten Winkel bei $L$ gegenüber.\\
(b) $\overline{LM}$ & 1\E & $\overline{LM} = \sqrt{8.5^2 - 4^2} = \sqrt{56.25} = 7.5\,\text{cm}$\\
(c) Fläche & 1 & Die Katheten stehen senkrecht aufeinander: $A = \tfrac{1}{2} \cdot 4 \cdot 7.5 = 15\,\text{cm}^2$. Folgepunkt: mit dem eigenen $\overline{LM}$.\\
\end{raster}
\fehler{Quadrate addiert ($\sqrt{8.5^2 + 4^2} \approx 9.39$, länger als die Hypotenuse): (b) 0; (c) als Folgepunkt ($\approx 18.79$) $\to$ 2 von 3\,P, sofern (a) stimmt.
Fläche als $\tfrac{1}{2} \cdot 4 \cdot 8.5$ (Hypotenuse als Höhe): (c) 0.}
\end{lpaufgabe}

\begin{lpaufgabe}{G6}{Giebelwand}{4 P}
\begin{raster}
(a) Ansatz & 1 & Die Höhe halbiert die Grundlinie: rechtwinkliges Dreieck mit den Katheten $4.8\,\text{m}$ und $3.6\,\text{m}$, die Schräge ist die Hypotenuse.\\
(a) Schräge & 1\E & $s = \sqrt{4.8^2 + 3.6^2} = \sqrt{36} = 6\,\text{m}$\\
(b) Fläche & 1\E & $A = \tfrac{1}{2} \cdot 9.6 \cdot 3.6 = 17.28\,\text{m}^2$\\
(c) Umfang & 1 & $U = 9.6 + 2 \cdot 6 = 21.6\,\text{m}$. Folgepunkt: mit der eigenen Schräge.\\
\end{raster}
\fehler{Ganze Breite statt halber ($\sqrt{9.6^2 + 3.6^2} \approx 10.25$): (a) beide Zeilen 0; (c) als Folgepunkt ($30.10\,\text{m}$) $\to$ 2 von 4\,P.
Fläche ohne $\tfrac{1}{2}$ ($34.56$): (b) 0.}
\end{lpaufgabe}

\begin{lpaufgabe}{G7}{Jonas' Rechnung}{2 P}
\begin{raster}
Fehler & 1 & Jonas addiert die Längen. Nach Pythagoras werden die Quadrate addiert, danach wird die Wurzel gezogen.\\
Richtig & 1\E & $c = \sqrt{2.5^2 + 6^2} = \sqrt{42.25} = 6.5\,\text{cm}$\\
\end{raster}
\fehler{Wurzel vergessen ($42.25$): Ergebniszeile 0. Nur «falsch» ohne Grund: Fehlerzeile 0.}
\end{lpaufgabe}

\needspace{14\baselineskip}
\section*{4 · Selbsteinschätzung}
\begin{tabularx}{\linewidth}{@{}l X@{}}\toprule
22 – 25 P & Die geprüften Teile sitzen. Wo du Punkte verloren hast: das Kapitel dieser Aufgabe nochmals (Zuordnung unten).\\
17 – 21 P & Das schwächste Kapitel nochmals: Tüfteln und Übungen des Kapitels, in dem du die meisten Punkte verloren hast.\\
11 – 16 P & Zurück zu den Kapiteln aller Aufgaben, in denen du Punkte verloren hast.\\
0 – 10 P & Zurück zu Kapitel 1 und von dort der Reihe nach weiter.\\\midrule
\multicolumn{2}{@{}p{\linewidth}@{}}{\small\textbf{Aufgabe $\to$ Kapitel:} G1 $\to$ Kapitel 1 (Winkel im Dreieck); G2, G3 $\to$ Kapitel 2 (Höhen, Halbierende, Mittelsenkrechte); G4 $\to$ Kapitel 3 (Fläche und Umfang), Seite $b$ mit Kapitel 4; G5, G6, G7 $\to$ Kapitel 4 (Pythagoras), Fläche und Umfang mit Kapitel 3}\\\bottomrule
\end{tabularx}
\end{document}
